
\subsection*{Remarque 3.8.}

Dans le cas \(X\) noethérien et \(Y\) régulier, 3.7 détermine entièrement le groupe abélien \(K^{\bullet}(X')\) à partir des \(\lambda\)-anneaux \(K^{\bullet}(X)\), \(K^{\bullet}(Y)\), des applications \(i_*,i^*\) entre ces anneaux, et de l'élément \(N\) de \(K^{\bullet}(Y)\) : en effet, \(K^{\bullet}(Y')\) peut se calculer en fonction de \(K^{\bullet}(Y)\) et de \(N\) d'après (VI 1.1). Nous verrons au paragraphe suivant que la structure multiplicative de \(K^{\bullet}(X')\) est déterminée par les mêmes données, de même que la \(\lambda\)-structure si l'on tensorise par \(\mathbb Q\) (cette restriction étant sans doute inutile : voir 4.5).

\section*{4. Immersions régulières et filtration du \(K^{\bullet}\)}

On garde dans ce paragraphe les notations et les hypothèses de 3.1.

\textbf{Lemme 4.1.} \uline{Supposons qu'il existe sur} \(Y\) \uline{un faisceau localement libre de type fini} \(\underline{N}'\), \uline{tel que}, \uline{si} \(N,N'\) \uline{sont les classes dans} \(K^{\bullet}(Y)\) \uline{de} \(\underline{N},\underline{N}'\), \uline{on ait} \(N=N'+2\). \uline{Alors la congruence suivante est vérifiée} :
\[
  \lambda_{-1}(F)\equiv 0 \ \mathrm{mod.}(1-L).
\]

Soit \(d\) le rang de \(\underline{N}\) (qu'on peut supposer constant) ; celui de \(\underline{F}\) est alors \(d-1\). Par définition :
\[
\begin{aligned}
\lambda_{-1}(F)
 &= \sum_{i\ge0}(-1)^i\lambda^i(F)\\
 &= (-1)^{d-1}\lambda^{d-1}(F-1)\\
 &= (-1)^{d-1}\lambda^{d-1}(F+L-1-L)\\
 &= (-1)^{d-1}\lambda^{d-1}(N-2+1-L).
\end{aligned}
\]
Or \(\lambda_t(1-L)=\lambda_t(1)/\lambda_t(L)=(1+t)/(1+Lt)\). Donc pour tout \(i\ge1\), \(\lambda^i(1-L)\) est divisible par \(1-L\). Par suite :
\[
\begin{aligned}
\lambda_{-1}(F)
  &\equiv (-1)^{d-1}\lambda^{d-1}(N-2) \qquad\mathrm{mod.}(1-L)\\
  &\equiv(-1)^{d-1}\lambda^{d-1}(N') \qquad\mathrm{mod.}(1-L).
\end{aligned}
\]
